\section{Outlook and Summary} \label{chapter:SUMMARY} \label{chapter:FUTURE_WORK}

We presented an approach to preserving access to previous states of objects in programming systems such as the Lively Kernel.
The approach is based on version-aware references that can be enabled for existing JavaScript environments without changing the \ac{vm}\@.
These references manage different versions of objects transparently, \ie they automatically resolve to one of multiple versions of an object; to which version in particular can easily be changed.
Thereby, different preserved states can be re-established.

Our approach uses a whole-system source transformation to ensure that for each object that is created, a proxy is created and returned instead of the object.
Thus, references to proxies are passed around and all access goes through the proxies.
Moreover, versions of an object are reclaimed together with their proxy by the ordinary garbage collector, keeping the memory overhead reasonable.

An issue with our current implementation is that the execution
overhead is not yet practical, with a slowdown of about three orders
of magnitude.  We assume that browser vendors have not spent
significant resources on optimizing the proxies, because the
ECMAScript 6 proxy specification is, at the time of this writing,
still being drafted.  It seems reasonable to assume that further
optimizations will be included in the browsers once the ECMAScript 6
specification is finalized.  Alternatively, instead of using proxies,
version-aware references could be implemented differently: using
source transformations the traps of our proxies could be inserted
directly around all relevant operations, removing one indirection.
Preliminary performance tests indicate that this alternative
implementation of version-aware references could be faster than
proxies as currently available in browsers. Using this technique for
re-implementing version-aware references, their overhead could
conceivably decrease to be within one order of magnitude.

Another issue is that, while our implementation allows to preserve and re-establish versions, these versions need to be created explicitly and there are no tools yet to find and manage versions.
Instead of putting the burden of remembering to create a new version on the programmer, the system could create versions of the runtime for any change to an object caused by an action of the programmer, as in \coexist.
Such actions would entail
\begin{inparaenum}[a)]
    \item manipulating properties of a morph directly with a halo tool or through drag and drop
    \item adding, removing, or editing a script of a morph or a method of a class
    \item evaluating a code snippet
    \item triggering code execution through a mouse or keyboard interaction.
\end{inparaenum}
This way, whenever programmers realize changes were inappropriate, they can undo their actions.

To complement this our system could support developers in
finding and re-establishing relevant states by storing and presenting
additional information about versions. Such information could include
timestamps, the kind of action that triggered preserving the version,
as well as which objects, classes, or files were changed, and how the
changes affected benchmarks and tests.

Provided a fast proxy implementation and proper tool support, the
proposed approach allows developers to focus on the programming task.
They can directly manipulate parts or evaluate code snippets to try
out ideas without having to anticipate errors. When errors do occur,
recovery is fast and easy to accomplish, because previous versions of
the system are immediately available.


%%% Local Variables:
%%% mode: latex
%%% TeX-master: "../main"
%%% End:
